\chapter{TYPE INFERENCE ALGORITHM}

\section{Overview}
% TODO: Explain what type inference means in this context.
% No explicit type annotations --- the compiler infers types
% from usage (literals, operations, assignments).

\section{Type Environment}
% TODO: Define the type environment (symbol table) \Gamma:
%   \Gamma : Identifier -> Type
% Explain how it is built up during parsing/checking.

\section{Inference Rules}

\subsection{Literals}
% TODO: Inference rules for integer, float, boolean, string literals.
%   e.g.  \Gamma |- 42 : Integer
%         \Gamma |- 3.14 : Float
%         \Gamma |- true : Boolean
%         \Gamma |- "hi" : String

\subsection{Arithmetic Expressions}
% TODO: Inference rules for +, -, *, /
% Cover both Integer and Float cases.
% Note: no overloading --- operands must have the same type.

\subsection{Boolean Expressions}
% TODO: Inference rules for == and !=
% Both sides must be arithmetic expressions of the same type.
% Result type is always Boolean.

\subsection{Assignment Statement}
% TODO: Inference rule for x = expr
% If x not in \Gamma: infer type from expr and add to \Gamma.
% If x in \Gamma: check that expr type matches stored type.

\subsection{If-Then-Else}
% TODO: Inference rule for if-then-else.
% Condition must be Boolean; both branches are type-checked independently.

\subsection{While-Loop}
% TODO: Inference rule for while-loop.
% Condition must be Boolean; body is type-checked.

\subsection{Function Definition and Call}
% TODO: Inference rules for function definition (parameter types inferred
% from body) and function call (argument types must match parameter types).

\section{Type Error Reporting}
% TODO: Describe how type errors are detected and reported to the user.
% e.g. type mismatch in arithmetic, incompatible comparison types, etc.

\section{Example Type Inference Traces}
% TODO: Walk through 2--3 examples showing step-by-step type inference.
% Example 1: simple arithmetic
% Example 2: if-then-else with boolean condition
% Example 3: function call
